% =============================================================================
% Capitulo 7 - Anexos
% =============================================================================
\chapter{Material complementario}
\label{cap:anexos}

\section{Acrónimos}
\label{anexo:acronimos}

\begin{acronym}[SigMF]
  \acro{TFG}{Trabajo Fin de Grado}
  \acro{UAV}{vehículo aéreo no tripulado (\emph{Unmanned Aerial Vehicle})}
  \acro{UAS}{sistema aéreo no tripulado (\emph{Unmanned Aircraft System})}
  \acro{DJI}{\emph{Da-Jiang Innovations}}
  \acro{RF}{radiofrecuencia}
  \acro{SDR}{radio definida por software (\emph{Software Defined Radio})}
  \acro{IQ}{componentes en fase y cuadratura}
  \acro{ISM}{banda industrial, científica y médica (\emph{Industrial, Scientific and Medical})}
  \acro{USRP}{\emph{Universal Software Radio Peripheral}}
  \acro{SigMF}{\emph{Signal Metadata Format}}
  \acro{STFT}{transformada de Fourier de tiempo corto (\emph{Short-Time Fourier Transform})}
  \acro{FFT}{transformada rápida de Fourier (\emph{Fast Fourier Transform})}
  \acro{FHSS}{salto de frecuencia (\emph{Frequency Hopping Spread Spectrum})}
  \acro{DSSS}{espectro ensanchado por secuencia directa (\emph{Direct Sequence Spread Spectrum})}
  \acro{OFDM}{multiplexación por división de frecuencias ortogonales (\emph{Orthogonal Frequency Division Multiplexing})}
  \acro{AFH}{salto de frecuencia adaptativo (\emph{Adaptive Frequency Hopping})}
  \acro{GFSK}{modulación por desplazamiento de frecuencia gaussiana (\emph{Gaussian Frequency Shift Keying})}
  \acro{AWGN}{ruido blanco gaussiano aditivo (\emph{Additive White Gaussian Noise})}
  \acro{SNR}{relación señal a ruido (\emph{Signal-to-Noise Ratio})}
  \acro{AGC}{control automático de ganancia (\emph{Automatic Gain Control})}
  \acro{CNN}{red neuronal convolucional (\emph{Convolutional Neural Network})}
  \acro{ResNet}{red residual (\emph{Residual Network})}
  \acro{PCA}{análisis de componentes principales (\emph{Principal Component Analysis})}
  \acro{YOLO}{\emph{You Only Look Once}}
  \acro{mAP}{precisión media promediada (\emph{mean Average Precision})}
  \acro{IoU}{intersección sobre la unión (\emph{Intersection over Union})}
  \acro{COCO}{\emph{Common Objects in Context}}
  \acro{GNSS}{sistema global de navegación por satélite (\emph{Global Navigation Satellite System})}
  \acro{NLOS}{sin línea de visión directa (\emph{Non Line of Sight})}
  \acro{RCS}{sección radar (\emph{Radar Cross Section})}
  \acro{GPU}{unidad de procesamiento gráfico (\emph{Graphics Processing Unit})}
\end{acronym}

\section{Convención de nombres de captura}
\label{anexo:nombres}

Para mantener trazabilidad se recomienda utilizar nombres de captura con campos explícitos:

\begin{lstlisting}[caption={Formato recomendado de nombre de captura.}]
<dron>_<label>_<distancia/nodrone>_<wifi_level>_<sesion>_<capNN>
\end{lstlisting}

Ejemplos:

\begin{lstlisting}[caption={Ejemplos de nombres de captura.}]
mini3_drone_d5_w1_s01_cap01
mini3_drone_d10_w3_s02_cap01
mini3_nodrone_w1_s01_cap03
mini3_nodrone_w3_s01_cap03
\end{lstlisting}

En el CSV de entrenamiento, la etiqueta de captura debe normalizarse como \texttt{drone} o \texttt{interference}, aunque el nombre mantenga el prefijo \texttt{nodrone} para indicar que el dron estaba apagado.

\section{Composición detallada del dataset v3}
\label{anexo:dataset_v3}

La \cref{tab:anexo_dataset_v3} detalla las \num{21}~capturas del dataset~v3 (\num{150}~espectrogramas cada una, \num{3150} en total), con su distancia, nivel de WiFi, frecuencia central y partición. El prefijo \texttt{mini3\_} se omite por brevedad. La partición se realiza por captura completa para evitar fuga de información entre entrenamiento, validación y prueba.

\begin{table}[htbp]
\centering
\footnotesize
\caption{Composición del dataset v3 por captura.}
\label{tab:anexo_dataset_v3}
\begin{tabular}{lccccc}
\toprule
\textbf{Captura} & \textbf{Clase} & \textbf{Dist. (m)} & \textbf{WiFi} & \textbf{$f_c$ (MHz)} & \textbf{Partición} \\
\midrule
drone\_d10\_w1\_s01\_cap01 & drone & 10 & 1 & 2450 & train \\
drone\_d10\_w1\_s01\_cap02 & drone & 10 & 1 & 2450 & train \\
drone\_d10\_w2\_s02\_cap01 & drone & 10 & 2 & 2460 & train \\
drone\_d10\_w2\_s02\_cap02 & drone & 10 & 2 & 2470 & val \\
drone\_d10\_w3\_s02\_cap01 & drone & 10 & 3 & 2460 & test \\
drone\_d10\_w3\_s02\_cap02 & drone & 10 & 3 & 2470 & train \\
drone\_d5\_w1\_s01\_cap01  & drone & 5  & 1 & 2450 & test \\
drone\_d5\_w1\_s01\_cap02  & drone & 5  & 1 & 2470 & val \\
drone\_d5\_w2\_s02\_cap01  & drone & 5  & 2 & 2450 & train \\
drone\_d5\_w2\_s02\_cap02  & drone & 5  & 2 & 2475 & train \\
drone\_d5\_w3\_s02\_cap01  & drone & 5  & 3 & 2450 & train \\
drone\_d5\_w3\_s02\_cap02  & drone & 5  & 3 & 2470 & train \\
nodrone\_w1\_s01\_cap01 & interference & n/a & 1 & 2422 & train \\
nodrone\_w1\_s01\_cap02 & interference & n/a & 1 & 2422 & train \\
nodrone\_w1\_s01\_cap03 & interference & n/a & 1 & 2422 & test \\
nodrone\_w2\_s01\_cap01 & interference & n/a & 2 & 2422 & train \\
nodrone\_w2\_s01\_cap02 & interference & n/a & 2 & 2422 & train \\
nodrone\_w2\_s01\_cap03 & interference & n/a & 2 & 2422 & train \\
nodrone\_w3\_s01\_cap01 & interference & n/a & 3 & 2422 & train \\
nodrone\_w3\_s01\_cap02 & interference & n/a & 3 & 2422 & val \\
nodrone\_w3\_s01\_cap03 & interference & n/a & 3 & 2422 & test \\
\midrule
\multicolumn{6}{l}{\textbf{Total:} train \num{2100} (14 capturas), val \num{450} (3 capturas), test \num{600} (4 capturas)} \\
\bottomrule
\end{tabular}
\end{table}

Las capturas de dron se registraron en frecuencias centrales de \SIrange{2450}{2475}{\mega\hertz} y las de interferencia a \SI{2422}{\mega\hertz}. Esta diferencia de portadora entre clases no afecta a la decisión del detector, que opera sobre la morfología local de las ráfagas tras la normalización frecuencial (\texttt{FrequencyNormalize}), pero se documenta aquí por transparencia; su eliminación mediante portadoras espejadas entre clases se contempla en el dataset~v4 propuesto en el \cref{cap:conclusiones}.

\section{Configuración de entrenamiento de los modelos}
\label{anexo:hiperparametros}

La \cref{tab:anexo_hiperparametros} resume la configuración de los tres modelos entrenados en el trabajo: el detector YOLOv8n de la primera etapa, el clasificador ResNet18 de seis clases de la segunda etapa y el clasificador binario ResNet18 del experimento P1.

\begin{table}[htbp]
\centering
\footnotesize
\caption{Hiperparámetros de los tres modelos entrenados.}
\label{tab:anexo_hiperparametros}
\begin{tabular}{lll}
\toprule
\textbf{YOLOv8n (Stage 1)} & \textbf{ResNet18 (Stage 2)} & \textbf{ResNet18 binario (P1)} \\
\midrule
Base: yolov8n.pt (COCO)      & Base: ResNet18 (ImageNet)   & Base: ResNet18 (ImageNet) \\
Entrada: 640$\times$640       & Entrada: 224$\times$224     & Entrada: 224$\times$224 \\
Lote: 16                      & Lote: 32                    & Lote: 32 \\
Épocas: 60 (baseline 15)      & Épocas: 30 (mejor en la 10) & Épocas: 20 (mejor en la 5) \\
Optimizador: SGD              & Optimizador: AdamW          & Optimizador: AdamW \\
Semilla: 42                   & lr $10^{-4}$, wd $10^{-4}$   & lr $10^{-4}$, wd $10^{-4}$ \\
Aumentos geom. desactivados   & CosineAnnealingLR           & CosineAnnealingLR \\
(fliplr, flipud, degrees,     & 3 grupos congelados         & 3 capas congeladas \\
shear, perspective = 0)       & FrequencyNormalize          & FrequencyNormalize \\
                              & Semilla: 42                 & Semilla: 42 \\
\bottomrule
\end{tabular}
\end{table}

\section{Parámetros del barrido de SNR}
\label{anexo:snr_params}

El barrido evalúa el clasificador de seis clases (\num{941}~segmentos de prueba) inyectando ruido en el dominio IQ a los niveles de SNR $\{+\infty,\, +20,\, +15,\, +10,\, +5,\, 0,\, -5,\, -10,\, -15,\, -20\}$~dB, regenerando el espectrograma con el mismo \emph{pipeline} de entrenamiento (STFT con ventana de Hann de \num{1024}~muestras, solape \SI{75}{\percent}, mapa de color parula, \num{1024}$\times$\num{576}~píxeles). Se comparan dos modelos de ruido calibrados a la misma potencia total:

\begin{itemize}
  \item \textbf{AWGN:} ruido gaussiano complejo $n[k]\sim\mathcal{CN}(0,\sigma^2)$, con $\sigma^2$ ajustada al SNR objetivo, distribuido uniformemente por todo el plano tiempo-frecuencia.
  \item \textbf{FHSS sintético:} \num{160}~ráfagas por ventana de \SI{0.1}{\second} (\num{1600}~saltos/s), cada una de \SI{366}{\micro\second} de duración y aproximadamente \SI{1}{\mega\hertz} de ancho, situadas en frecuencia e instante aleatorios dentro del \SI{90}{\percent} de la banda capturada, con modulación GFSK simplificada (chirp de $\pm$\SI{250}{\kilo\hertz}) y ventana de Hann en amplitud.
\end{itemize}

\section{Pseudocódigo del auto-etiquetado YOLO}
\label{anexo:pseudocodigo_autolabel}

\begin{lstlisting}[style=python,caption={Pseudocódigo de generación automática de cajas YOLO.}]
for image_path in spectrograms:
    img = read_image(image_path)
    intensity = rgb_to_intensity(img)
    threshold = percentile(intensity, 99.8)
    mask = intensity > threshold
    components = connected_components(mask)

    for comp in components:
        box = bounding_box(comp)
        area = box.width * box.height

        if area < 8:
            continue
        if area > 0.005 * image_width * image_height:
            continue

        class_id = class_from_capture_name(image_path)
        write_yolo_box(image_path, class_id, box)
\end{lstlisting}

\section{Pseudocódigo del detector YOLO}
\label{anexo:pseudocodigo_detector_yolo}

\begin{lstlisting}[style=python,caption={Pseudocódigo de inferencia del detector YOLO.}]
iq = read_iq_window(capture, duration=0.1)
spectrogram = make_spectrogram(iq)
detections = yolo.predict(spectrogram)

for det in detections:
    if det.confidence < threshold:
        continue
    if det.class_name == "drone":
        alert = True
        crop = crop_region(spectrogram, det.box)
\end{lstlisting}

\section{Pseudocódigo de búsqueda y confirmación}
\label{anexo:pseudocodigo_busqueda}

Una versión operativa del sistema podría recorrer varias frecuencias, generar espectrogramas cortos y bloquear la frecuencia cuando el detector confirme presencia de dron:

\begin{lstlisting}[style=python,caption={Pseudocódigo de búsqueda y confirmación con YOLO.}]
centers = [2.410e9, 2.425e9, 2.440e9, 2.455e9, 2.470e9, 2.485e9]

for fc in centers:
    tune_usrp(fc)
    iq = capture_short(duration=0.25)
    spec = make_spectrogram(iq)
    detections = yolo.predict(spec)

    if has_drone_detection(detections, conf=0.80):
        iq2 = capture_short(duration=0.25)
        spec2 = make_spectrogram(iq2)
        detections2 = yolo.predict(spec2)

        if has_drone_detection(detections2, conf=0.80):
            lock_frequency(fc)
            capture_long(duration=3.0)
            break
\end{lstlisting}

\section{Pseudocódigo de identificación de modelos}
\label{anexo:pseudocodigo_modelos}

\begin{lstlisting}[style=python,caption={Pseudocódigo de clasificación multiclase como segunda etapa.}]
spec = make_spectrogram(iq_window)
detections = yolo.predict(spec)

for det in detections:
    if det.class_name != "drone":
        continue
    crop = crop_region(spec, det.box)
    probs = model_multiclass.predict(crop)
    model_id = argmax(probs)
    confidence = max(probs)
\end{lstlisting}

\section{Disponibilidad de código, modelos y datos}
\label{anexo:disponibilidad}

El código de adquisición, generación de espectrogramas, auto-etiquetado, entrenamiento y evaluación, la fuente \LaTeX{} de esta memoria y los resultados consolidados (métricas y figuras) están disponibles en el repositorio del proyecto: \url{https://github.com/edumateoss/spectre-main}. Los \emph{datasets} de espectrogramas, las capturas IQ en formato SigMF y los pesos de los modelos entrenados (detector YOLOv8n y clasificador ResNet18) no se incluyen en el repositorio por su tamaño y se entregan como ficheros asociados a esta memoria. La estructura del repositorio, los comandos de reproducción y las versiones de las dependencias se documentan en el fichero \texttt{README} incluido.
